\subsection{李群的正则元}

定义 2. G 的一个元素称为正则的，如果它关于 G 的伴随表示是正则的。

换句话说（命题 1），元素 \(g \in \mathbf{G}\) 是正则的，如果对所有元素 \(g'\)（\(g\) 在 \(\mathbf{G}\) 中的一个邻域内的元素），\(\operatorname{Ad}(g') - 1\) 的幂零空间的维数等于 \(\operatorname{Ad}(g) - 1\) 的幂零空间的维数。

命题 2. 设 G' 是 k 上的一个有限维李群，f 是 G 到 G' 的一个开态射。G 的正则元素在 f 下的像是 G' 的正则元素。若 f 的核包含于 G 的中心，则元素 \(g \in G\) 正则当且仅当 \(f(g)\) 正则。

实际上，设 \(\mathfrak{g}'\) 是 \(G'\) 的李代数，\(\mathfrak{h}\) 是 \(\mathfrak{g}\) 中由 \(Tf|\mathfrak{g}\) 的核给出的理想。设 \(\rho\) 是 \(G\) 在 \(\mathfrak{h}\) 上的线性表示，由 \(\rho(g) = \operatorname{Ad} g|\mathfrak{h}\)（对所有 \(g \in G\)）定义，并设 \(\operatorname{Ad} \circ f\) 是 \(G\) 在 \(\mathfrak{g}'\) 上的线性表示，由 \(f\) 与 \(G'\) 的伴随表示复合而成。这些线性表示定义了 \(G\)-模的一个正合列：

\[
0 \to \mathfrak {h} \to \mathfrak {g} \to \mathfrak {g} ^ {\prime} \to 0.
\]

由引理 2，\(r_{\mathrm{Ad}} = r_{\rho} + r_{\mathrm{Ad}\circ f}\)。由于 \(r_{\mathrm{Ad}\circ f} = r_{\mathrm{Ad}}\circ f\) 且 \(f\) 是开映射，有 \(r_{\mathrm{Ad}\circ f}^{0} = r_{\mathrm{Ad}}^{0}\circ f\)。因而：

\[
r _ {\mathrm{Ad}} - r _ {\mathrm{Ad}} ^ {0} = r _ {\rho} - r _ {\rho} ^ {0} + \left(r _ {\mathrm{Ad}} - r _ {\mathrm{Ad}} ^ {0}\right) \circ f.
\]

这样，若 g 正则，则 \((r_{\mathrm{Ad}} - r_{\mathrm{Ad}}^{0})(f(g)) = 0\)，这表示 \(f(g)\) 正则。若 f 的核包含于 G 的中心，

\[
r _ {\rho} (g) = r _ {\rho} ^ {0} (g) = \dim \mathfrak {h}
\]

对所有 \(g \in G\) 成立。因而，若 \(f(g)\) 正则，则 \(r_{\mathrm{Ad}}(g) = r_{\mathrm{Ad}}^{0}(g)\)，换言之，g 正则。

命题 3. 设 \(G_{1}\) 与 \(G_{2}\) 是 k 上的两个有限维李群。元素 \((g_{1}, g_{2})\)（\(G_{1} \times G_{2}\) 的）是正则的，当且仅当 \(g_{1}\) 与 \(g_{2}\) 分别是 \(G_{1}\) 与 \(G_{2}\) 的正则元素。

由命题 2，条件是必要的。我们证明它是充分的。对所有 \(g = (g_{1}, g_{2}) \in \mathrm{G}_{1} \times \mathrm{G}_{2}\)，\(r_{\mathrm{Ad}}(g) = r_{\mathrm{Ad}}(g_{1}) + r_{\mathrm{Ad}}(g_{2})\)。由引理 1 (ii) 推出 \(r_{\mathrm{Ad}}^{0}(g) = r_{\mathrm{Ad}}^{0}(g_{1}) + r_{\mathrm{Ad}}^{0}(g_{2})\)。若 \(g_{1}\) 与 \(g_{2}\) 正则，则 \(r_{\mathrm{Ad}}^{0}(g_{1}) = r_{\mathrm{Ad}}(g_{1})\) 且 \(r_{\mathrm{Ad}}^{0}(g_{2}) = r_{\mathrm{Ad}}(g_{2})\)，于是 \(r_{\mathrm{Ad}}^{0}(g) = r_{\mathrm{Ad}}(g)\)，这表示 \(g\) 正则。

引理 3. 设 \(a \in \mathbf{G}\)，并设 \(\mathfrak{m}\) 是 \(\mathfrak{g}^1(a)\) 在 \(\mathfrak{g}\) 中的一个补。设 \(\mathbf{U}\) 是 \(\mathfrak{g}\) 中 0 的一个邻域，exp 是 \(\mathbf{U}\) 到 \(\mathbf{G}\) 的一个指数映射。映射

\[
f: (x, y) \mapsto (\exp y) a (\exp x) (\exp y) ^ {- 1}
\]

（从 \((\mathfrak{g}^1 (a)\times \mathfrak{m})\cap \mathrm{U}\) 到 G）在 \((0,0)\) 处是平展的。

映射 \(x \mapsto a(\exp x)\) 与 \(y \mapsto (\exp y)a(\exp y)^{-1}\) 在 0 处的切线性映射分别是映射 \(x \mapsto ax\) 与 \(y \mapsto ya - ay = a(a^{-1}ya - y)\)（从 \(\mathfrak{g}\) 到 \(T_aG = ag\)；第三章，§3，第 12 节，命题 46）。因此，\(f\) 在 \((0,0)\) 处的切映射是映射 \((x,y) \mapsto ax + a(a^{-1}ya - y) = a(x + a^{-1}ya - y)\)（从 \(\mathfrak{g}^1(a) \times \mathfrak{m}\) 到 \(ag\)）。这个映射是单射。事实上，若 \(x \in \mathfrak{g}^1(a), y \in \mathfrak{m}\) 且 \(x + a^{-1}ya - y = 0\)，则 \((\mathrm{Ad}(a) - 1)y = \mathrm{Ad}(a)x \in \mathfrak{g}^1(a)\)，因为 \(\mathrm{Ad}(a)\mathfrak{g}^1(a) \subset \mathfrak{g}^1(a)\)。这蕴含 \(y \in \mathfrak{g}^1(a)\)，从而 \(y = 0\)。由于 \(\mathrm{Ad}(a)\) 在 \(\mathfrak{g}^1(a)\) 上是单射，推出 \(x = 0\)。由于 \(\dim \mathfrak{g} = \dim \mathfrak{g}^1(a) + \dim \mathfrak{m}\)，这表明 \(f\) 在 \((0,0)\) 处是平展的。

命题 4. 设 \(a \in G\)，\(H\) 是 \(G\) 的一个李子群芽，其李代数为 \(\mathfrak{g}^1(a)\)。映射 \((b, c) \mapsto cabc^{-1}\)（从 \(H \times G\) 到 \(G\)）在 \((e, e)\) 处是一个浸没。

实际上，设 \(\mathfrak{m}\) 是 \(\mathfrak{g}^1(a)\) 在 \(\mathfrak{g}\) 中的一个补，exp 是 \(G\) 的一个指数映射，定义在开邻域 \(U\)（\(0\) 在 \(\mathfrak{g}\) 中的一个开邻域）上。我们可以选取 \(U\)，使得 \(\exp(U \cap \mathfrak{g}^1(a)) \subset H\)。映射 \(f: (x, y) \mapsto (\exp x, \exp y)\) 是 \((0, 0)\) 在 \(\mathfrak{g}^1(a) \times \mathfrak{m}\) 中的一个邻域上的解析映射，取值于 \(H \times G\)。由引理 3，\(f\) 与映射 \(\varphi: (b, c) \mapsto cabc^{-1}\) 的复合在 \((0, 0)\) 处是平展的。由此推出 \(\varphi\) 在 \(f(0, 0) = (e, e)\) 处是一个浸没。

命题 5. 设 \(a \in G\)，并设 \(W\) 是 \(e\) 在 \(G\) 中的一个邻域。则存在一个邻域 \(V\)（\(a\) 的邻域），具有下列性质：对所有 \(a' \in V\)，存在元素 \(g \in W\)，使得 \(\mathfrak{g}^1(a') \subset \mathrm{Ad}(g)\mathfrak{g}^1(a)\)。

置 \(\mathfrak{g}^1 = \mathfrak{g}^1(a)\)，并设 \(\mathfrak{g} = \mathfrak{g}^1 + \mathfrak{g}^+\) 是 \(\operatorname{Ad}(a) - 1\) 的 Fitting 分解（§1，第 1 节）。设 H 是 G 的一个李子群芽，其李代数为 \(\mathfrak{g}^1\)。对所有 \(h \in \mathrm{H}\)，\(\operatorname{Ad}(h)\mathfrak{g}^1 \subset \mathfrak{g}^1\)。由于 \([\mathfrak{g}^1, \mathfrak{g}^+] \subset \mathfrak{g}^+\)，存在 \(e\) 在 H 中的一个邻域 U，使得 \(\operatorname{Ad}(h)\mathfrak{g}^+ \subset \mathfrak{g}^+\) 对所有 \(h \in \mathrm{H}\) 成立。由于 \(\operatorname{Ad}(a)-1\) 在 \(g^{+}\) 上的限制是双射，可以选取 U，使得 \(\operatorname{Ad}(ah)-1\) 在 \(g^{+}\) 上的限制对所有 \(h\in U\) 是双射。于是 \(\mathfrak{g}^{1}(ah)\subset\mathfrak{g}^{1}(a)=\mathfrak{g}^{1}\) 对所有 \(h\in U\) 成立。由命题 4，\(\operatorname{Int}(\mathrm{W})(a\mathrm{U})\) 是 a 在 G 中的一个邻域。若 \(a'\in\operatorname{Int}(\mathrm{W})(a\mathrm{U})\)，则 \(a'=g(ah)g^{-1}\)，其中 \(g\in W\) 且 \(h\in U\)；由此推出 \(\mathfrak{g}^{1}(a')=\operatorname{Ad}(g)\mathfrak{g}^{1}(ah)\subset\operatorname{Ad}(g)\mathfrak{g}^{1}(a)\)。

推论. 设 \(G^{*}\) 是 G 的一个开子群。若 \(a \in G\) 正则，则存在 a 的一个邻域 V，使得对所有 \(a' \in V\)，\(\mathfrak{g}^{1}(a')\) 与 \(\mathfrak{g}^{1}(a)\) 在 \(\mathrm{Ad}(\mathrm{G}^{*})\) 下共轭。

